\chapter{Introduction}
\label{ch:intro}

\section{Background and Context}
Question-answer based story generation represents a paradigm shift in
the field of Generative Artificial Intelligence, transforming how
narratives are created and consumed. Unlike traditional storytelling
systems that present users with predetermined content, Q\&A-based
systems position the user as an \emph{active participant} in the
creative process. The interaction is no longer a single prompt-response
exchange but a sustained dialogue from which a story emerges.

The evolution of artificial intelligence has witnessed remarkable
progress in natural-language understanding and generation. Early
rule-based systems gave way to statistical models, which were
subsequently superseded by neural-network approaches. The advent of
the transformer architecture \cite{vaswani2017attention} and
attention mechanisms enabled Large Language Models (LLMs) to achieve
unprecedented performance across language tasks, including creative
writing and narrative generation
\cite{brown2020gpt3,openai2023gpt4,touvron2023llama,jiang2023mistral}.

Modern LLMs such as OpenAI's GPT-3/GPT-4 series \cite{openai2023gpt4},
Google's BERT and T5 \cite{devlin2019bert}, Meta's LLaMA family
\cite{touvron2023llama}, and BART \cite{lewis2020bart} have demonstrated
remarkable capabilities in understanding natural-language queries,
capturing subtle user intent, maintaining contextual coherence across
extended interactions, and generating text that exhibits creativity,
emotional depth and stylistic versatility. Trained on vast corpora
comprising trillions of tokens, these models have internalised
patterns of human language use that allow them to produce narratives
often indistinguishable from human-written content in many contexts.

\section{Motivation and Significance}
The demand for interactive storytelling spans several important
application domains:

\begin{itemize}
    \item \textbf{Educational technology} -- adaptive story-based
    tutorials that respond to learner understanding, style and
    interest; narrative-based learning consistently improves
    retention and engagement.
    \item \textbf{Entertainment and gaming} -- interactive fiction
    and procedurally adaptive role-playing where the user actively
    shapes outcomes through choices and responses.
    \item \textbf{Therapeutic applications} -- narrative therapy,
    trauma processing and anxiety-management protocols, in which
    personalised stories scaffold emotional reflection.
    \item \textbf{Content creation and marketing} -- scalable
    generation of personalised marketing narratives, children's
    stories, brand storytelling and creative writing assistance.
    \item \textbf{Accessibility} -- making creative storytelling
    accessible to individuals who lack traditional writing skills
    but possess rich imagination expressible through conversation.
\end{itemize}

By integrating LLMs with a guided question-answer framework, this
project aims to enhance user engagement while ensuring story
coherence, creativity, relevance and emotional resonance. The
objective is to design a robust, scalable system capable of
generating high-quality narratives from user-provided responses while
maintaining consistency across extended interactions.

\section{Scope and Boundaries of the Study}
This dissertation develops a comprehensive Q\&A-based interactive
storytelling system using generative AI technologies. The scope
encompasses:

\begin{itemize}
    \item Designing structured input mechanisms that capture diverse
    user preferences across four narrative dimensions (setting,
    character, theme, plot).
    \item Implementing sophisticated context-extraction algorithms
    that identify narrative elements from conversational input.
    \item Integrating state-of-the-art LLM-based story generation
    with custom prompt-engineering strategies.
    \item Developing intuitive output-delivery systems with feedback
    and continuation mechanisms.
    \item Augmenting the textual narrative with three orchestrated
    modalities --- per-segment scene illustration, voice narration
    and ambient music --- conditioned on detected genre and
    segment content.
    \item Surfacing reader-facing analytics through a Story Map and
    a Character Relationship Graph.
    \item Providing an end-user authoring layer for custom answers
    and entirely user-added Q\&A pairs.
\end{itemize}

The system is designed to be practical for real-world deployment,
scalable to handle concurrent users, and suitable for adaptation
across multiple domains.

\section{Research Questions}
This project addresses several key research questions:
\begin{itemize}
    \item How can structured Q\&A interactions effectively capture
    user preferences for narrative generation, including preferences
    for genre, tone, character, plot complexity and themes?
    \item What architectural approaches enable LLMs to maintain
    narrative coherence and consistency across extended multi-turn
    interactions while incorporating continuous user input?
    \item How can context extraction and management be optimised to
    preserve relevant story elements while allowing natural narrative
    evolution based on user direction?
    \item What evaluation metrics effectively measure the quality of
    interactive storytelling systems across coherence,
    personalisation, engagement and creativity?
    \item How can the system balance user agency in directing
    narratives against the need for coherent, well-structured stories
    that follow established narrative conventions?
\end{itemize}

\section{Contributions of this Dissertation}
The principal contributions of this work are:

\begin{enumerate}
    \item A \textbf{novel Q\&A framework} for interactive storytelling
    that enables real-time personalisation through conversational
    interaction across four canonical narrative dimensions.
    \item A \textbf{seven-module backend pipeline} that cleanly
    separates context extraction, prompt engineering and arc control
    while remaining LLM-provider agnostic.
    \item A \textbf{multimodal augmentation layer} comprising
    sentence-scored text-to-image synthesis, Web-Speech-API narration
    and genre-conditioned ambient music streaming, with a serial
    throttling queue to prevent rate-limit cascades.
    \item Two \textbf{D3-based reader analytics} (Story Map branch
    tree and Character Relationship Graph) with an explicit
    co-appearance disclaimer.
    \item An \textbf{end-user authoring layer} that admits free-form
    custom answers per question and entirely user-added Q\&A pairs,
    routed through a dedicated \texttt{custom\_elements} context
    channel into the LLM prompt.
    \item A reproducible \textbf{evaluation methodology} with
    quantitative latency, quality and consistency measurements over
    50 stories and 15 user-study participants.
\end{enumerate}

\section{Structure of the Report}
The remainder of this report is organised as follows.
Chapter~\ref{ch:litreview} surveys related work in symbolic, neural
and interactive storytelling and crystallises the research gap.
Chapter~\ref{ch:problem} formalises the problem statement, lists
objectives and scopes the work.
Chapter~\ref{ch:schedule} presents the project schedule as a Gantt
chart spanning August 2025 -- May 2026.
Chapter~\ref{ch:proposed} describes the proposed system covering
feasibility study, framework, hardware/software details, design and
methodology.
Chapter~\ref{ch:impl} documents implementation details module by
module with annotated UI snapshots, including the rendered Q\&A
interface, generated story segments, the Story Map and the Character
Relationship Graph.
